\section{Roadmap and Scope Discipline}
\label{sec:roadmap}

This section documents what is in scope for v1.x, what is scheduled for v1.4/v2, and what is explicitly out of scope.

\subsection{Per-project release state (as of 2026-05-21)}

\begin{center}
\small
\begin{tabularx}{\textwidth}{@{}lllX@{}}
\toprule
\textbf{Project} & \textbf{Version} & \textbf{Date} & \textbf{Major milestones} \\
\midrule
\tesseraid{}    & v1.1 & 2026-05-21 & Eight-language byte-identical implementations; HLC variant; cross-project variant-byte rejection contract pinned. \\
\tesserahash{}  & v1.0 & 2026-05-21 & Eight-language byte-identical implementations; SHAKE-256 tail; cross-rejection of cross-project bytes. \\
\tesseraseal{} & v1.3 (+hybrid) & 2026-06-13 & Eight-language alpha for suites \code{0x01}/\code{0x02}; hybrid \code{0x06} as a 4-way conformance contract (Python, Rust, Go, C) plus library-level hybrid in C++, TypeScript, CPython ext; plugin protocol + \texttt{.tseal-key-v1.1} keyfile; CLI binary. \\
\bottomrule
\end{tabularx}
\end{center}

\subsection{Open backlog (v1.4 / v2)}

\textbf{TesseraSeal hybrid \code{0x06} port status (as of 2026-06-13):}

\begin{enumerate}
\item TSEAL-00030: Hybrid suite \code{0x06} port to C (OpenSSL 3.5+ \texttt{EVP\_PKEY\_X25519}). \textbf{Shipped} --- promotes hybrid to a 4-way byte-identical conformance contract (Python, Rust, Go, C) and closes the largest single gap.
\item TSEAL-00031: Hybrid \code{0x06} C++ wrapper (routes to the C entry points). \textbf{Shipped} --- library-level, verified by in-suite \texttt{hybrid\_*.json} KAT byte-identity (no \tesseraseal{} conformance runner yet).
\item TSEAL-00032: Hybrid \code{0x06} in TypeScript (\texttt{@noble/curves} X25519, pinned 2.0.1). \textbf{Shipped} --- library-level, verified by in-suite KAT byte-identity.
\item TSEAL-00034: Hybrid \code{0x06} in the CPython extension (wraps the C reference). \textbf{Shipped} --- library-level, verified by in-suite KAT byte-identity.
\item TSEAL-00033: Hybrid \code{0x06} in Java (SunJCE \texttt{XDH} keypair). \textbf{Open} --- the sole remaining hybrid gap, pending a JDK 24+ build/verify environment.
\end{enumerate}

\textbf{Other TesseraSeal backlog:}

\begin{enumerate}
\item Add \tesseraseal{} conformance runners for C++, TypeScript, Java, and the CPython extension so their (already library-level) suite coverage --- including hybrid \code{0x06} where implemented --- is harness-enforced rather than in-suite-test-enforced.
\item Rust AES-NI/ARMv8 intrinsics for 1\,MiB AEAD throughput. Currently regressed because \texttt{aes-gcm 0.10} does not enable ARMv8 intrinsics on stable Rust without a nightly cfg.
\item Suite \code{0x07} (HQC contingency) is reserved on the registry (TSEAL-00048) but \emph{not} allocated; v1.x decoders MUST reject it. Allocation is gated on NIST HQC FIPS finalization or a Module-LWE confidence loss.
\end{enumerate}

\textbf{TesseraID v1.x:} no open backlog. The format is stable; future work is operational (benchmarks, additional reference implementations, documentation polish).

\textbf{TesseraHash v1.x:} CPython extension is the largest single gap; mirrors the Python API.

\textbf{Cross-project:} sibling project \texttt{pqsign} (variant byte \code{0x05}? or further into the reserved-future range) is an explicit \emph{out-of-v1.x} candidate. The minimal scope would be ML-DSA \citep{fips204} signatures of a payload-with-context, using a \tesseraid{} or \tesserahash{} as the canonical payload. Not committed; documented as a possibility so the reserved-future range is interpretable.

\subsection{Scope discipline}

\begin{itemize}
\item \textbf{No new variant bytes without a whitepaper update.} \code{0x00} and \code{0xFF} are permanently reserved; \code{0x05}--\code{0xFE} are reserved-future. Any new allocation requires (a) a normative whitepaper update in the owning project, (b) conformance-vector additions to every existing project verifying cross-rejection, and (c) sign-off as a security-relevant change.
\item \textbf{No third-party crypto dependencies without justification.} Each project's whitepaper pins per-language sources for its primitives. New dependencies require explicit justification in the PR description.
\item \textbf{No unsafe Rust, no cgo, no FFI shortcuts.} Implementations stay pure within their language. The CPython extensions are C-against-CPython-API, not FFI bridges.
\item \textbf{No format-version mutations of an existing variant.} A variant byte \emph{is} a version. A change to the layout, the security argument, or the algorithm choices under an existing byte is a spec break. New algorithms get new bytes --- and the same rule applies one level down in \tesseraseal{}'s \emph{suite}-byte namespace at header offset 7: the hybrid construction took the new suite byte \code{0x06} inside variant \code{0x04} rather than mutating suite \code{0x01}.
\end{itemize}

\subsection{What \emph{would} trigger a v2 of the whole toolkit}

\begin{itemize}
\item \textbf{A structural break of Module-LWE.} Would force a successor envelope format: a hypothetical \tesseraseal{} v2 takes a new \emph{variant} byte allocated from the reserved-future range \code{0x05}--\code{0xFE}. (Note the namespace distinction: \code{0x06} is a \emph{suite} byte inside variant \code{0x04}'s suite namespace at header offset 7, not a variant byte --- the two one-byte namespaces are disjoint, so the variant-byte allocation table in Section~\ref{sec:shared-namespace} is unaffected by suite-byte assignments.) \tesseraid{} and \tesserahash{} are unaffected by ML-KEM breaks because they make no PQ-cryptography claim.
\item \textbf{A serious break of SHAKE-256.} Would force \tesserahash{} to a new variant byte with a different hash function (KMAC256? a stream cipher in PRF mode? a different sponge?). \tesseraseal{}'s KDF schedule would also need revision because HKDF-SHA3-256 depends on SHAKE-256. \tesseraid{} is unaffected.
\item \textbf{A serious break of AES-256-GCM.} Would force \tesseraseal{} to a new suite under a new byte using a different AEAD. \tesseraid{} and \tesserahash{} are unaffected.
\item \textbf{A change to the variant-byte allocation discipline.} Would force every project to a new wire-format version, indicated by a new variant byte. This is the worst-case scenario and is exactly what the immutable-variant-byte rule is designed to prevent.
\end{itemize}

None of these is anticipated in the v1.x lifetime. The roadmap is conservative on purpose: \tessera{} is intended to outlive any individual deployment, which means most of its versions look like ``a new entry was added to the variant table; everything else is unchanged.''
